\documentclass[10pt,a4paper]{amsart}
\usepackage[margin=19mm]{geometry}
\usepackage{amsmath,amssymb,mathtools}
\usepackage[scr=rsfs,cal=euler]{mathalfa}
\usepackage{xfrac}
\usepackage[unicode,hidelinks,bookmarks=false]{hyperref}
\newcommand{\ie}{i.e.\ }
\newcommand{\cf}{cf.\ }
\newcommand{\resp}{respectively\ }
\newcommand{\Subsection}{\subsection}
\providecommand{\nobreakdash}{\nobreak\hbox{-}}
\setlength{\emergencystretch}{2em}
\multlinegap0pt
\usepackage{amssymb}
\usepackage{mathtools}
\usepackage[shortlabels]{enumitem}
\newtheorem{lemma}{Lemma}
\title{Homological $k$-systole in $n$-manifolds with positive intermediate curvature}
\author{Jingche Chen, Han Hong}
\date{Source: arXiv:2601.12461v2 (CC BY 4.0)}
\begin{document}
\maketitle

\noindent\emph{Source and licence.} Homological $k$-systole in $n$-manifolds with positive intermediate curvature. Version arXiv:2601.12461v2 (CC BY 4.0), distributed by arXiv under the Creative Commons Attribution 4.0 International licence, http://creativecommons.org/licenses/by/4.0/, as recorded in the arXiv licence metadata for this exact version and shown on the abstract page https://arxiv.org/abs/2601.12461v2. Full licence notice: \texttt{licenses/arxiv-2601.12461-CC-BY-4.0.txt}.\par\smallskip

\noindent\textit{Editorial scope.} Counted unit: the article's lemma inherit (source0.tex lines 964--973). The article's complete original proof of it is delivered with it (source0.tex lines 974--994, one \texttt{proof} environment of 21 lines). No mathematics is written, repaired or re-derived: the statement and the proof are the article's own lines, in the article's own notation, and no retained line is substituted or re-flowed.  No further source-local context is needed: every cross-reference of every retained line resolves inside the two delivered spans.  Nothing was omitted from any retained span, so the package carries no omission record.  No retained line contains a citation, so the wrapper carries no bibliography and no bibliography entry.  No retained line contains a figure environment, an image-inclusion command or drawing code, so no image is delivered and the package declares no asset dependency.  Editorial formatting only, in addition: an AMS \texttt{amsart} preamble that restates the article's own declarations and macros used by the retained lines, namely the environments \texttt{lemma} on separate counters, together with the standard packages those lines need.  The retained environments print in this document as Lemma 1 in order of appearance, whereas the article prints the same items with the numbering of its own full document.  The complete native source of the article as distributed in the arXiv e-print for this exact version is bundled unchanged as \texttt{sources/192952/source0.tex}.
\begin{lemma}\label{topology inherit}
    Let $N$ be an $n$-dimensional closed orientable manifold. Suppose $3\leq n\leq 7$, $1\leq m\leq n-1$ and there exists a smooth map $f:N\rightarrow E\times
      \mathbb{T}^{m-1}$ with nonzero degree. 
      Let $0\leq k\leq n-2$ and let $W\subset N$ be a closed, oriented
      $(n-k)$-dimensional smooth submanifold representing the homology class
      $[N]\frown(f^{*}(d\theta_{1})\smile \cdots\smile f^{*}(d\theta_{k}))\in
      H_{n-k}(N)$. Then there exists an area-minimizing hypersurface
      $\Sigma\subset W$ such that $\Sigma$ represents the homology class
      $[N]\frown(f^{*}(d\theta_{1})\smile \cdots\smile f^{*}(d\theta_{k+1}))\in H_{n-k-1}(N)$.
\end{lemma}

\begin{proof}
    Let $H_{k+1}\subset N$ be a smooth closed oriented hypersurface representing
      $[N]\frown f^{*}(d\theta_{k+1})$. After a small perturbation, we may assume that
      $H_{k+1}$ intersects $W$ transversely. Then $H_{k+1}\cap W$ is a smooth closed oriented
      hypersurface in $W$. We use $P$ to denote the Poincare duality map. By the standard
      intersection formula for Poincare duals,
    \[
    \begin{aligned}
    &[H_{k+1}\cap W] = P(P^{-1}([H_{k+1}])\smile P^{-1}([W]))\\
    &= [N]\frown(f^{*}(d\theta_{1})\smile \cdots\smile f^{*}(d\theta_{k+1}))\neq 0\in H_{n-k-1}(N).
    \end{aligned}
    \]
    Then
    \[
    [H_{k+1}\cap W]\neq 0 \in H_{n-k-1}(W).
    \]
    Thus, there exists $\Sigma\subset W$ which is a minimizing hypersurface representing $[H_{k+1}\cap W]\in H_{n-k-1}(W)$. Moreover,
    \[
    [\Sigma]= [H_{k+1}\cap W]\in H_{n-k-1}(N).
    \]
\end{proof}

\end{document}
